\section{Finite observable modules produce complete rational point lists}
\label{sec:finite-observable-modules}

The earlier finite-algebra interface proves that a multiplication-closed
spanning certificate of size $D$ bounds the number of rational zeros by
$D$.  There is more information in the same certificate.  Its matrices
also bound every coordinate and every reduced denominator.  The resulting
finite search returns the complete rational zero set, with no externally
supplied height bound.  This is a strengthening of that interface, not a
claim that an arbitrary cardinality bound determines a search cutoff.
The additional information is the supplied multiplication action.

The central observation needs neither a determinant nor a linearly
independent basis.  A rational eigenvalue of an integer matrix is an
integer: after scaling an eigenvector to a primitive integer vector,
the denominator would otherwise divide all its coordinates.  We give
the finite proof because it both simplifies the argument and identifies
exactly the arithmetic needed in iPRA.

\subsection{Primitive integer eigenvectors}

All finite vectors and matrices in this section have at least one
coordinate.  A vector $w\in\mathbb Z^D$ is \emph{primitive} when the
greatest common divisor of $|w_1|,\ldots,|w_D|$ is one.  Rational
numbers are represented by reduced numerator--positive-denominator
pairs.  Matrix multiplication uses row vectors: $(wA)_j=
\sum_{i=1}^D w_i A_{ij}$.

\begin{lemma}[Primitive scaling]
\label{lem:primitive-scaling}
From a supplied nonzero vector $v\in\mathbb Q^D$ one can primitively
recursively construct a nonzero rational $c$ and a primitive integer
vector $w$ such that $w=cv$.  If $vM=\lambda v$, then
$wM=\lambda w$.
\end{lemma}

\begin{proof}
Write the entries of $v$ with positive denominators, and let $L$ be
their product.  Then $u=Lv$ is an integer vector and is nonzero.
Compute $g=\gcd(|u_1|,\ldots,|u_D|)$ by the Euclidean algorithm,
folded along the finite list.  Since $u$ is nonzero, $g>0$.
Every $u_i$ is divisible by $g$, so $w_i=u_i/g$ is an integer and
the gcd of its absolute coordinates is one.  Take $c=L/g$.
Multiplying the displayed eigenvector equation by $c$ gives the final
assertion.  Products, gcd, exact integer division with its failure
branch, and finite folds are primitive recursive.  Extended gcd
additionally produces integers $e_i$ satisfying
$\sum_i e_iw_i=1$, if that form of primitivity is preferred.
\end{proof}

\begin{lemma}[Integer eigenvalue and an explicit bound]
\label{lem:integer-eigenvalue}
Let $A\in\mathbb Z^{D\times D}$, let $w\in\mathbb Z^D$ be
primitive, and let $\lambda\in\mathbb Q$.  If
$wA=\lambda w$, then $\lambda\in\mathbb Z$.  Moreover, with
\[
 R(A)=\max_{1\le j\le D}\sum_{i=1}^D|A_{ij}|,
\]
one has $|\lambda|\le R(A)$.
\end{lemma}

\begin{proof}
Write $\lambda=a/b$ in reduced form, where $b>0$ and
$\gcd(|a|,b)=1$.  For each coordinate $j$, multiplication by $b$
in the eigenvector equation gives
\[
 a w_j=b\sum_iw_iA_{ij}.
\]
Thus $b$ divides $a w_j$, and coprimality gives $b\mid w_j$.
One elementary justification uses a Bezout identity $ra+sb=1$:
multiply it by $w_j$ and substitute the preceding divisibility.
Since $w$ is primitive, there are integers $e_j$ with
$\sum_j e_jw_j=1$.  Therefore $b\mid1$.  Positivity of $b$
forces $b=1$.

For the bound, choose the first index $j$ attaining
$W=\max_i|w_i|$.  The finite maximum and its first maximizing
index are obtained by a bounded scan.  Primitivity implies $W>0$.
The $j$-th coordinate gives
\[
 |\lambda|W
   =\left|\sum_iw_iA_{ij}\right|
   \le\sum_i|w_i||A_{ij}|
   \le W\sum_i|A_{ij}|
   \le WR(A).
\]
Cancel the positive integer $W$.  No matrix inverse, determinant,
characteristic polynomial, or spectral theorem was used.
\end{proof}

\begin{corollary}[A rational eigenvalue lies on a finite grid]
\label{cor:eigenvalue-grid}
Let $M\in\mathbb Q^{D\times D}$ and let $\Delta>0$ be an
integer such that $A=\Delta M$ has integer entries.  If a nonzero
$v\in\mathbb Q^D$ satisfies $vM=xv$, then
\[
 x=\frac z\Delta
 \quad\text{for some }z\in\mathbb Z\text{ with }|z|\le R(A).
\]
In particular, if $x=a/q$ is reduced with $q>0$, then
$q\mid\Delta$, $q\le\Delta$, and $|a|\le R(A)$.
\end{corollary}

\begin{proof}
Primitive scaling gives $wM=xw$ with $w$ primitive, and hence
$wA=(\Delta x)w$.  Lemma~\ref{lem:integer-eigenvalue} says that
$z=\Delta x$ is an integer with the asserted bound.  To reduce
$z/\Delta$, divide numerator and denominator by
$\gcd(|z|,\Delta)$.  This proves the denominator divisibility
and numerator bounds, including $z=0$.
\end{proof}

\subsection{The observable-module principle}

This principle applies before choosing a particular polynomial geometry.
Let $S$ be a class of finitely coded objects with a decidable validity
predicate.  Let $u_1,\ldots,u_r:S\to\mathbb Q$ be the rational
observables whose possible values are wanted.  Fix rational-valued
functions $b_1,\ldots,b_D$, rational numbers $c_1,\ldots,c_D$,
and rational matrices $M_1,\ldots,M_r$.  Suppose there are proved
local contracts, for every valid $p$, of the form
\begin{align}
  1&=\sum_{i=1}^D c_i b_i(p),
      \label{eq:observable-unit}\\
  u_k(p)b_j(p)&=\sum_{i=1}^D (M_k)_{ij}b_i(p)
      &&(1\le k\le r, 1\le j\le D).
      \label{eq:observable-action}
\end{align}
The names ``observable module'' and ``spanning family'' refer only to
these displayed finite identities.  One need not introduce an abstract
module or assert that the functions $b_i$ are linearly independent.

\begin{theorem}[Finite observable image]
\label{thm:finite-observable-image}
Choose positive common denominators $\Delta_k$ for the matrices,
put $A_k=\Delta_kM_k\in\mathbb Z^{D\times D}$, and put
$R_k=R(A_k)$.  Every valid $p$ satisfies
\[
 (u_1(p),\ldots,u_r(p))\in
 \prod_{k=1}^r\left\{\frac z{\Delta_k}:z\in\mathbb Z,
                                      -R_k\le z\le R_k\right\}.
 \label{eq:observable-grid}
\]
Consequently the image has at most $\prod_k(2R_k+1)$ elements.
If all functions and the validity test have specified primitive
recursive implementations and the local contracts have iPRA proofs,
the displayed implication also has an iPRA proof, uniformly in the
finite matrix and denominator data.
\end{theorem}

\begin{proof}
For a fixed valid $p$, let $v=(b_1(p),\ldots,b_D(p))$.
The unit identity implies that $v\ne0$: if all entries vanished,
their indicated linear combination would be zero, contradicting one.
The action identity says precisely that $vM_k=u_k(p)v$.
Apply Corollary~\ref{cor:eigenvalue-grid} separately to every $k$.
The product of the finite coordinate grids contains every image value.
Each grid has exactly $2R_k+1$ entries because $\Delta_k>0$,
and finite tuple enumeration gives the cardinality bound.

For the iPRA assertion, the arbitrary coded input $p$ is a parameter
of the proof.  All remaining quantifiers in the matrix-vector
identities, gcd folds, Bezout checks, and maximum searches range over
supplied finite lists.  Compile them to Boolean folds.  The Euclidean
and list identities just proved then use quantifier-free induction.
There is no quantification over all functions or all possible programs.
\end{proof}

If observables fail to determine $p$, this theorem supplies a finite
\emph{image}, not a list of all input objects.  An infinite set can
have a constant observable.  Complete enumeration therefore needs
actual rational coordinates, or a separately proved finite reconstruction
procedure from the observable values.  This is a useful distinction:
the theorem can also classify finite invariant values of an infinite
family without pretending to enumerate that family.

\subsection{Polynomial certificates and the enumeration theorem}

Let $F_1,\ldots,F_s\in\mathbb Q[X_1,\ldots,X_r]$ be finitely
encoded polynomials.  A \emph{unit-and-action certificate} comprises
polynomials $b_1,\ldots,b_D$, a rational vector $c$, rational
matrices $M_k$, and coefficient-polynomial witnesses for the identities
\begin{align}
 1-\sum_i c_i b_i&=\sum_{\ell=1}^s h_\ell F_\ell,
                \label{eq:module-unit-certificate}\\
 X_kb_j-\sum_i (M_k)_{ij}b_i
       &=\sum_{\ell=1}^s h_{kj\ell}F_\ell.
                \label{eq:module-action-certificate}
\end{align}
The checker expands each side and compares the finitely many
coefficients.  The earlier spanning-certificate interface is the
special case $b_1=1$, $c_1=1$, and $c_i=0$ for $i>1$.
Our relaxed unit condition is occasionally convenient when a generating
family is supplied without a distinguished constant member.

\begin{theorem}[Complete rational enumeration from a finite certificate]
\label{thm:module-enumeration}
From a valid unit-and-action certificate, a fixed primitive recursive
algorithm produces a duplicate-free list $L$ satisfying
\[
  p\in L\quad\Longleftrightarrow\quad
           F_1(p)=\cdots=F_s(p)=0
       \qquad(p\in\mathbb Q^r).
\]
It examines at most $\prod_{k=1}^r(2R_k+1)$ candidate tuples,
where $\Delta_k$, $A_k$, and $R_k$ are as in
Theorem~\ref{thm:finite-observable-image}.  Every coordinate
$p_k=a_k/q_k$ of a rational zero, in reduced form, obeys
\[
 q_k\mid\Delta_k,
 \qquad |a_k|\le R_k,
 \qquad \max(|a_k|,q_k)\le\max(R_k,\Delta_k).
\]
No linear independence of the $b_i$ and no pairwise commutation of
the matrices is required by the checker or the proof.
\end{theorem}

\begin{proof}
At a common zero $p$, substitute its coordinates into every checked
polynomial identity.  The right sides vanish, leaving precisely
\eqref{eq:observable-unit} and \eqref{eq:observable-action} with
$u_k(p)=p_k$.  The finite observable image theorem yields the full
coordinate grid and the reduced-height bounds.

Now enumerate all integer tuples $(z_1,\ldots,z_r)$ with
$-R_k\le z_k\le R_k$, replace them by
$(z_1/\Delta_1,\ldots,z_r/\Delta_r)$ in canonical rational form,
and retain exactly those tuples for which each $F_\ell$ evaluates
to zero.  Every retained tuple is a common zero by the retention test.
Every common zero is visited by the grid bound and retained, proving
completeness.  Distinct integer tuples give distinct rational tuples
because each denominator $\Delta_k$ is fixed and nonzero.  Thus no
additional mathematical selection or deduplication principle is needed.

To implement the Cartesian product with uniform arity, put
$B_k=2R_k+1$ and $B=\prod_kB_k$.  For each $0\le n<B$, recover
its mixed-radix digits $d_k<B_k$ by successive remainder and quotient
operations, and set $z_k=d_k-R_k$.  These arithmetic operations and
the finite decoding loop are PR.  Polynomial evaluation, rational
normalization, list filtering, and the certificate check are also
fixed PR algorithms.  On an invalid certificate the total program
returns a failure flag; its soundness and completeness assertion is
conditional on the checker accepting.

Only the individual eigenvector equations were used.  Matrices that
fail to commute on unused linear directions can therefore still be
valid certificates, and redundant generators cause no difficulty.
\end{proof}

The theorem does \emph{not} find a certificate for an arbitrary input
system.  Certificate production must come from a supplied algebraic
construction, a separately bounded elimination algorithm, or an
independently checked finite output.  A general curve cannot possess
such a certificate for all of its coordinates when it has infinitely
many rational points.  The theorem does not turn a finite rational-point
cardinality assertion into a height bound for that curve.

\paragraph{A stronger algebraic obstruction for a whole plane curve.}
In fact, finite rational-point cardinality need not make the certificate
exist at all.  Let $F\in\mathbb Q[X,Y]$ be monic of positive degree
$m$ in $Y$.  Then $\mathbb Q[X]$ injects into
$\mathbb Q[X,Y]/(F)$.  To prove this without any geometric theorem,
suppose $q(X)=H(X,Y)F(X,Y)$.  If $H\ne0$, its leading
$Y$-coefficient survives multiplication by the monic $F$, so the
right side has $Y$-degree $\deg_YH+m\ge1$; a nonzero $q(X)$
has $Y$-degree zero.  Therefore $H=0$ and $q=0$.

A unit-and-action certificate of size $D$ for the sole equation
$F=0$ would span the entire quotient over $\mathbb Q$: use the
unit vector $c$ for the class of $1$, then the action matrices
to reduce each successive variable multiplication, and add the
reductions for a polynomial's finitely many monomials.  In particular,
$1,X,\ldots,X^D$ would reduce to $D+1$ vectors of length $D$.
Finite rational elimination gives a nontrivial dependence between
them, yielding a nonzero $q(X)\in(F)$ and contradicting the preceding
paragraph.  Thus no such full-coordinate certificate exists, even
when the curve has only finitely many rational points.  For example,
$Y^2+X^2+1=0$ has no rational points, by the order properties of
$\mathbb Q$, yet its monic plane-curve quotient still has this
obstruction.  This polynomial-identity certificate interface handles
finite algebraic sections and systems with supplied finite spanning
data; rational finiteness
of a positive-dimensional algebraic curve is a different kind of
information.  All parts of this obstruction can be read finitely:
the contradiction uses only $D+1$ coefficient vectors, one rational
dependence, and one polynomial degree comparison.

The earlier size-$D$ zero-count theorem remains available.  Its
interpolation argument extends to the relaxed unit condition by using
$c$ as the starting vector for the reduction of the constant polynomial.
Consequently the returned list has at most $D$ entries, although the
coarser grid can be larger.  The grid is a search bound; the algebraic
dimension bound is a stronger count bound.  They serve different roles.

\subsection{A constructor from separated annihilator identities}
\label{sec:separated-annihilators}

The certificate interface becomes an algorithm rather than a mere input
specification whenever one has separated polynomial annihilators with
checked ideal-membership witnesses.  Suppose, for each $1\le k\le r$,
the input supplies a monic univariate polynomial
\[
 q_k(T)=T^{d_k}+\sum_{i=0}^{d_k-1}a_{ki}T^i,
 \qquad d_k\ge1,
\]
and multivariate polynomials $h_{k\ell}$ satisfying
\[
 q_k(X_k)=\sum_{\ell=1}^s h_{k\ell}(X)F_\ell(X).
 \label{eq:separated-annihilator}
\]
These are exact finite polynomial identities.  Merely knowing that
$q_k$ vanishes at some observed solutions is insufficient.

\begin{theorem}[Constructing the spanning certificate]
\label{thm:annihilator-to-module}
From the identities \eqref{eq:separated-annihilator}, one fixed
primitive recursive constructor produces a valid unit-and-action
certificate of size
\[
                         D=\prod_{k=1}^r d_k.
\]
Consequently it produces a complete rational-point list for the
original system.  If $\Delta_k$ clears the denominators of
$a_{k0},\ldots,a_{k,d_k-1}$, a valid coordinate grid uses
\[
 R_k=\max\left(\Delta_k,
               \sum_{i<d_k}|\Delta_k a_{ki}|\right).
 \label{eq:annihilator-grid-radius}
\]
\end{theorem}

\begin{proof}
Index the spanning family by the exponent box
\[
 E=\{e=(e_1,\ldots,e_r):0\le e_k<d_k\},
 \qquad b_e=X_1^{e_1}\cdots X_r^{e_r}.
\]
Use a mixed-radix enumeration of $E$, with the all-zero exponent vector
first.  The unit certificate is $b_0=1$.

Fix $k$ and $e$.  If $e_k<d_k-1$, multiplication by $X_k$
simply sends $b_e$ to the next basis monomial $b_{e+\varepsilon_k}$.
Its matrix column has one entry one and all other entries zero, and
the action residual is identically zero.

If $e_k=d_k-1$, put $\widehat b_e=\prod_{j\ne k}X_j^{e_j}$.
Then
\[
 X_kb_e
   -\left(-\sum_{i<d_k}a_{ki}b_{e[k\leftarrow i]}\right)
   =\widehat b_e q_k(X_k)
   =\sum_{\ell=1}^s
       (\widehat b_e h_{k\ell})F_\ell.
\]
Thus the column of $M_k$ has entry $-a_{ki}$ at the distinct rows
$e[k\leftarrow i]$, with zero elsewhere; the displayed products
$\widehat b_eh_{k\ell}$ are its explicit residual witnesses.
No repeated reduction or termination theorem is needed.  There is
one case distinction and one polynomial multiplication per boundary
column.

After multiplication by $\Delta_k$, a non-boundary column has
absolute sum $\Delta_k$, and a boundary column has absolute sum
$\sum_{i<d_k}|\Delta_ka_{ki}|$.  Hence
\eqref{eq:annihilator-grid-radius} bounds every column sum.
When $d_k=1$ there are no non-boundary columns, so the first term
of the maximum can be omitted; retaining it is harmless.
Apply Theorem~\ref{thm:module-enumeration}.

All exponent vectors, matrix entries, and residual witness polynomials
are produced by finite loops over $E$, the coordinate list, and the
input coefficient lists.  Mixed-radix decoding, multiplication of
polynomials, and the product $D$ are PR uniformly in their finite
input codes.  The correctness invariant is precisely the displayed
action identity for every already generated column.  Its quantifiers
are bounded by the constructed finite dimensions, so Booleanization
gives the required iPRA induction formula.
\end{proof}

The input membership witnesses carry the real elimination information.
The theorem does not assert that such witnesses can be found by an
unbounded search.  It does prove a uniform bounded constructor from
them.  If a supplied identity already has the form
$1=\sum_\ell h_\ell F_\ell$, the correct special case returns the
empty list immediately: substitution at a hypothetical common zero
would give $1=0$.

\paragraph{Optional converse through one adjugate identity.}
The two algebraic interfaces are equivalent as certificate formats,
although the direct coordinate-grid proof needs no determinant.
Indeed, suppose a unit-and-action certificate is given.  Fix $k$,
put $C= X_k I-M_k$, and put $q_k(T)=\det(TI-M_k)$.
The action identities say, for the row $b=(b_1,\ldots,b_D)$,
\[
 bC=r,\qquad r_j=\sum_\ell h_{kj\ell}F_\ell.
\]
Multiply by $\operatorname{adj}C$ and use the finite cofactor
identity $C\operatorname{adj}C=(\det C)I$.  Coordinate $j$ gives
\[
 q_k(X_k)b_j
   =\sum_\ell
       \left(\sum_t h_{kt\ell}(\operatorname{adj}C)_{tj}\right)
       F_\ell.
\]
Multiply by $c_j$, sum over $j$, and use the unit identity
$1-\sum_jc_jb_j=\sum_\ell h_\ell F_\ell$.  The result is an
explicit identity $q_k(X_k)=\sum_\ell u_{k\ell}F_\ell$, with
\[
 u_{k\ell}
   =q_k(X_k)h_\ell+
     \sum_j c_j\sum_t h_{kt\ell}(\operatorname{adj}C)_{tj}.
\]
Here $q_k$ is monic of degree $D$.  Finite permutation sums compute
the determinant and cofactors, so this is also a PR certificate
conversion; the usual cofactor cancellation proves its correctness.
The converse is useful when another application expects separate
elimination polynomials.  It is not needed for denominator bounds,
and it may produce a much larger certificate than the original one.

\subsection{Worked example: a finite section of an elliptic curve}

Consider the two integer polynomials
\[
 F_1=4y^2-4x^3+x-4,
 \qquad F_2=2x^2-x.
\]
The first equation alone is the affine cubic
$y^2=x^3-x/4+1$.  The second cuts out a finite section.  Use the
four generators
\[
             (b_1,b_2,b_3,b_4)=(1,x,y,xy).
\]
This basis and its action witnesses can be constructed automatically
from the two separated identities
\[
 x^2-x/2=F_2/2,
 \qquad y^2-1=\bigl(F_1+(2x+1)F_2\bigr)/4.
\]
Theorem~\ref{thm:annihilator-to-module} uses degrees $(2,2)$
and therefore returns precisely four monomial generators.
The unit identity is $b_1=1$.  In the row-vector convention above,
the multiplication matrices are
\[
 M_x=\begin{pmatrix}
 0&0&0&0\\
 1&1/2&0&0\\
 0&0&0&0\\
 0&0&1&1/2
 \end{pmatrix},
 \qquad
 M_y=\begin{pmatrix}
 0&0&1&0\\
 0&0&0&1\\
 1&0&0&0\\
 0&1&0&0
 \end{pmatrix}.
\]
For completeness, all nonzero residual identities are
\begin{align*}
 x b_2-\tfrac12 b_2&=\tfrac12 F_2,\\
 x b_4-\tfrac12 b_4&=\tfrac12 yF_2,\\
 y b_3-b_1&=\tfrac14 F_1+\tfrac14(2x+1)F_2,\\
 y b_4-b_2&=\tfrac14 xF_1+\tfrac14x(2x+1)F_2.
\end{align*}
The other four action residuals are identically zero.  Expanding the
third identity gives $y^2-1$ on both sides, and multiplying it by
$x$ gives the fourth.  Thus the certificate is a collection of eight
short polynomial identities, not an assumption about a quotient field.

Take $\Delta_x=2$ and $\Delta_y=1$.  The absolute column sums of
$A_x=2M_x$ are $(2,1,2,1)$, and those of $A_y=M_y$ are all one.
Hence $R_x=2$ and $R_y=1$.  Every rational common zero lies among
the fifteen candidate points
\[
 x\in\{-1,-1/2,0,1/2,1\},
 \qquad y\in\{-1,0,1\}.
\]
Filtering the two equations gives exactly
\[
                 (0,1),\ (0,-1),\ (1/2,1),\ (1/2,-1).
\]
One can verify the last filter symbolically as well: $F_2=0$
implies $x=0$ or $x=1/2$ by the integral-domain law; at either
value $F_1=0$ reduces to $y^2=1$, and therefore $y=1$ or $y=-1$.
The point count four attains the size of the spanning family.  The
denominator information is genuinely part of the result: the argument
forces $x$ to have reduced denominator dividing two and forces $y$
to be an integer before any solution is selected.

\subsection{Principal-open conditions and finite chart covers}

A finite algebra certificate can also treat equations together with a
polynomial inequation $G(x)\ne0$.  Introduce one extra rational
coordinate $t$ and the equation
\[
                           tG(x)-1=0.
\]
Every solution of the original system with $G(x)\ne0$ lifts uniquely
to the extended system by $t=1/G(x)$.  Conversely, an extended solution
forces $G(x)\ne0$ and projects to an original solution.  A valid
unit-and-action certificate for the extended system therefore gives
a complete finite list for that principal-open part.  Projection may
be followed by a bounded exact deduplication pass, although uniqueness
of $t$ already prevents repetitions coming from a single chart.

For projective equations, use the disjoint charts in which coordinate
$j$ is the first nonzero coordinate.  Normalize it to one, set all
earlier coordinates to zero, and substitute into the homogeneous
equations.  A supplied finite algebra certificate for every resulting
affine system gives a complete projective rational-point list: each
projective point belongs to exactly one such chart and has exactly one
normalized representative there.  Thus several small algebraic
certificates can replace one large elimination, without an appeal to
compactness.  This construction is conditional on the actual chart
certificates; no claim of general rational-point decidability is hidden
in the chart decomposition.

\subsection{The exact iPRA and Lean obligations}

The complete human proof reduces to finite arithmetic.  A practical
formalization should use the following dependency order.
\begin{enumerate}
\item Define signed integer and reduced rational codes, finite vectors,
matrix multiplication, mixed-radix tuple decoding, and sparse polynomial
evaluation.  Invalid codes receive specified default values.  Prove
normalization and Euclidean arithmetic by quantifier-free induction.
\item Prove a signed finite-vector Bezout lemma: if the gcd of the
absolute coordinates is one, the extended-gcd fold returns a coefficient
vector $e$ with $e\cdot w=1$.  Track the current gcd and the coefficient
identity as a bounded invariant.
\item Prove primitive scaling, then the integer-eigenvalue lemma using
the displayed coordinate divisibility identities.  The bound uses a
finite maximum, the triangle inequality for integer sums, and cancellation
of a positive integer.  No determinant library is required.
\item Define the certificate checker by coefficient equality for
\eqref{eq:module-unit-certificate} and
\eqref{eq:module-action-certificate}.  Prove its substitution lemma:
acceptance plus vanishing of the input equations implies the evaluated
unit and action identities.
\item Prove grid completeness and the finite filter specification.
The grid invariant states that each processed mixed-radix index has
been decoded correctly and retained exactly when all equation tests
succeed.  The input rational point is a free parameter when proving
that its grid index is visited.
\item Give explicit PR syntax for these fixed algorithms, compile the
bounded invariants to quantifier-free formulas, and build actual
\texttt{Prf PRAAx} derivations.  Only then apply the repository's
existing translation theorem into the native scalar presentation.
\end{enumerate}

The accompanying \path{lean/FiniteAlgebraTargets.lean} supplies
fully defined sparse polynomial operations, the certificate checker,
the Cartesian-grid producer, and the exact enumeration target.
Section~\ref{sec:formalization} gives the complete dependency order,
distinguishing checked host lemmas from the final object derivation.

The general observable theorem is a schema for fixed PR observables
with supplied local proofs.  The polynomial specialization is uniform
over arbitrary finite polynomial and certificate codes because its
evaluator executes one fixed polynomial grammar.  These two kinds of
uniformity must not be confused with a total universal evaluator for
all primitive recursive functions.
